\documentclass[11pt]{article}
\usepackage{amsmath,amssymb,geometry,hyperref}
\geometry{margin=1in}
\newcommand{\D}[1]{|D_{17}^{(#1)}\rangle}
\newcommand{\Tr}{\operatorname{Tr}}
\begin{document}
\title{A radical family of optimal third-order AC planes for seventeen qubits}
\author{Catalogue calculation}
\date{30 October 2026}
\maketitle
We give a two-real-parameter analytical family of two-dimensional third-order
anticoherent (AC) subspaces. A simple rational parameter choice gives a compact
exact witness; a second, exactly specified rational choice reconstructs the
reported numerical plane in \texttt{ACLiege.jl} to its numerical accuracy.
No equality between an exact state and rounded source decimals is asserted.

Throughout, $\D r$ is a normalized Dicke state with $r$ excitations.
All square roots used below are the nonnegative real square roots, and the
stated domain conditions make every denominator nonzero.

\section{The analytical construction}
Choose real parameters $p,r$, and put $q=241/768-p$. Define the two probability
lists, indexed by $i=0,\ldots,4$,
\begin{align*}
P={}&(p-65/768,\;799/768-4p,\;6p-731/768,\;255/256-4p,\;p),\\
Q={}&(q+65/768,\;505/768-4q,\;6q-341/768,\;539/768-4q,\;q).
\end{align*}
Next define
\begin{align*}
X&=-\frac{(1024p-255)(3072p-425)}{2895532523520},\\
Y&=-\frac{p(768p-241)}{13056},\\
Z&=\frac{25690112p^2-9670144p+770185}{68941250560},\\
z&=\sqrt X,\qquad w=\frac{Z+i\sqrt{XY-Z^2}}{\sqrt X},
\end{align*}
and
\[
\alpha=(-1820,3,-14/33,1,0),\quad
\beta=(-3,2/455,-1/2145,0,1),\quad x_i=\alpha_i z+\beta_i w.
\]
The letters $z,w$ here are auxiliary scalars, not sphere coordinates.
For $i=0,1,2,3$ set
\[
h=(14,2002,2002,14),\qquad
\ell_i^2=\frac{h_i^2P_{i+1}Q_i}{\binom{17}{4i+4}\binom{17}{4i+1}}.
\]
The second family parameter gives a unit complex number
\[
\eta=\frac{1-r^2+2ir}{1+r^2},\qquad
S=\ell_0+\ell_1\eta,\qquad R=|S|^2,
\]
\[
d=\frac{R+\ell_2^2-\ell_3^2}{2},\qquad
\Delta=\ell_2^2R-d^2.
\]
Four closed-quadrilateral edges are
\[
e_0=\ell_0,\quad e_1=\ell_1\eta,\quad
 e_2=-\frac{S(d+i\sqrt\Delta)}{R},\quad e_3=-S-e_2.
\]
The opposite sign of $i\sqrt\Delta$ is also allowed. We use the plus sign
for both parameter choices below.
Finally, the finite recurrence
\[
a_0=\sqrt{P_0},\qquad
a_{i+1}=\frac{a_i e_i}{h_i x_i}\quad(0\le i\le3),\qquad
b_i=\overline{\frac{x_i}{a_i}}
\]
provides every coefficient explicitly using rational operations and radicals.
The two basis vectors are
\begin{equation}\label{basis}
|u\rangle=\sum_{i=0}^4\sqrt{\binom{17}{4i}}a_i\D{4i},\qquad
|v\rangle=\sum_{i=0}^4\sqrt{\binom{17}{4i+1}}b_i\D{4i+1}.
\end{equation}
This is a finite radical expression, not an implicit polynomial-root definition.
The domain is
\[
P_i,Q_i>0,\qquad X>0,\qquad XY-Z^2>0,\qquad R>0,\qquad\Delta\ge0.
\]

\section{A simple exact point}
Take $p=5/24$ and $r=1$, so $\eta=i$. The probability lists simplify to
\[
P=\frac1{768}(95,159,229,125,160),\qquad
Q=\frac1{768}(146,181,145,215,81).
\]
The phase constants are
\[
X=\frac{5375}{1737319514112},\qquad
Y=\frac{45}{34816},\qquad Z=-\frac{232931}{124094251008}.
\]
The two relevant discriminants are strictly positive:
\[
XY-Z^2=\frac{1046084177}{2199911876176644145152},\qquad
\Delta=\frac{1057861}{11350123846041600}.
\]
Thus every coefficient in \eqref{basis} is a well-defined radical. The
machine-readable \texttt{basis.json} expands the finite recurrence into
ordinary symbolic expressions for all thirty-six Dicke components.

\section{Proof that the entire span is third-order AC}
The lists $P$ and $Q$ were obtained by solving the four diagonal moment
equations on supports $\{0,4,8,12,16\}$ and $\{1,5,9,13,17\}$. Directly,
\[
\sum_{i=0}^4P_i(4i)^a=\sum_{i=0}^4Q_i(4i+1)^a
=\frac1{18}\sum_{j=0}^{17}j^a,\qquad a=0,1,2,3.
\]
Occupied indices within each vector differ by at least four, so each
three-qubit self reduction is diagonal. These moment identities make
each of its four diagonal entries equal to $1/4$.

To verify both cross reductions, write the binomial-scaled coefficients
as in \eqref{basis}. The three possible entries with index difference one
are proportional to the rows of
\[
Kx=0,\qquad
K=\begin{pmatrix}
1&1001&3003&91&0\\
0&1092&10296&1092&0\\
0&91&3003&1001&1
\end{pmatrix}.
\]
Both $\alpha$ and $\beta$ lie in $\ker K$, hence the identity holds.
The probability constraints on these products are exactly
\[
|x_i|^2=\alpha_i^2X+\beta_i^2Y+2\alpha_i\beta_iZ
=\frac{P_iQ_i}{\binom{17}{4i}\binom{17}{4i+1}}.
\]
They are rational identities in $p$. Conversely, compatibility of this
kernel with the two moment families forces $p+q=241/768$.

The quadrilateral construction gives
\[
\sum_{i=0}^3e_i=0,\qquad |e_i|^2=\ell_i^2.
\]
For example, $|e_2|^2=(d^2+\Delta)/R=\ell_2^2$ and
$|e_3|^2=R+\ell_2^2-2d=\ell_3^2$.
The recurrence therefore implies inductively
\[
|a_i|^2=\frac{P_i}{\binom{17}{4i}},\qquad
|b_i|^2=\frac{Q_i}{\binom{17}{4i+1}}.
\]
The fourth possible cross entry, with index difference three, vanishes because
\[
\sum_{i=0}^3 h_i a_{i+1}\overline{b_i}
=\sum_{i=0}^3h_i\frac{a_{i+1}}{a_i}x_i
=\sum_{i=0}^3e_i=0.
\]
These are every potentially nonzero cross entry. The two supports are disjoint,
and each probability list sums to one, so the basis is orthonormal. Consequently
\[
\boxed{R_3(|u_a\rangle\langle u_b|)=\delta_{ab}I_4/4,
\qquad (u_1,u_2)=(u,v).}
\]
Every normalized complex superposition in the span is therefore third-order AC.
The primary verifier independently checks all $4\cdot4^2=64$ ordered scalar
cross-reduction entries at the simple point, using normalized Dicke splitting.

At that point the two displayed basis vectors have exactly AC order three:
their fourth-reduction first diagonal entries are respectively
\[
(\rho_4^u)_{00}=1099/5440=1/5+11/5440,\qquad
(\rho_4^v)_{00}=653/3264=1/5+1/16320.
\]
This does not classify the AC order of every superposition in the plane.

\section{Analytical reconstruction of the reported plane}
For the numerical example \texttt{\(\Pi\)\_17\_2\_3} at line 939 of
\texttt{ACLiege.jl}, use the exact rational parameters
\[
p_* =\frac{18370114016312017}{10^{17}},\qquad
r_* =\frac{47980997065595143}{10^{17}}.
\]
These are another member of the same radical family. The verifier
\texttt{verify\_family.py} checks all polynomial family identities and proves
that this rational point satisfies the domain inequalities using rational
square-root enclosures. Its numerical reduction to the source convention uses
one common collective $z$ rotation and independent basis phases. At eighty-digit
working precision, the maximum coefficient discrepancies from the supplied
rounded vectors are $1.73\times10^{-10}$ and $2.48\times10^{-10}$.
The exact existence argument is the family proof, not this numerical agreement.
The simple point $p=5/24,r=1$ is not claimed to be the original numerical plane.

\section{Optimality of the dimension}
For any third-order AC subspace of dimension $k$, take Hilbert--Schmidt
orthonormal rectangular irreducible spin-tensor maps
\[
A_{\ell m}:\operatorname{Sym}^{17}\mathbb C^2\to
\operatorname{Sym}^{16}\mathbb C^2,\qquad \ell=1/2,3/2.
\]
There are six maps. Every product $A_{\ell m}^\dagger A_{\ell'm'}$ has
tensor rank at most three, so the polarized AC conditions give
\[
\langle A_{\ell m}u_a,A_{\ell'm'}u_b\rangle
=\frac{\delta_{ab}\delta_{\ell\ell'}\delta_{mm'}}{18}.
\]
Thus $6k$ nonzero orthogonal vectors must fit in the seventeen-dimensional
output space: $6k\le17$, hence $k\le2$. The explicit family attains this bound:
\[
\boxed{K(17,3)=2.}
\]
The upper bound is unrestricted: no sparsity, symmetry or realness assumption
is imposed on competing subspaces. No classification of all optimal planes
or claim of literature priority is made.
\end{document}
